0586: \paragraph{读前自检：可逆马尔可夫链。}可逆链满足细致平衡$\rho_i a_{ij}=\rho_j a_{ji}$；请对$i$求和并用第$j$行和为1，核对$\sum_i\rho_i a_{ij}=\rho_j$，从而细致平衡推出平稳性。\par
0587: 
0588: \knowledgeanchor{SLM-19-06-001}{可逆马尔可夫链}
0589: \index{可逆马尔可夫链}
0590: \begin{definition}[细致平衡与可逆性]
0591: 若概率分布$\rho$满足
0592: \begin{equation}\label{eq:V5-C01-detailed-balance}
0593: \rho_i a_{ij}=\rho_j a_{ji},\qquad i,j\in\mathcal S,
0594: \end{equation}
0595: 则称$A$关于$\rho$满足细致平衡，链关于$\rho$可逆。等式两边都是平稳权重下沿同一条边相反方向的一步概率流。
0596: \end{definition}
0597: 本讲义只使用式\eqref{eq:V5-C01-detailed-balance}或等价的双向前向条件概率表达，不使用把时间反向条件概率与前向条件概率混为同一个式子的写法。
0598: 
0599: % CORE-DERIVATION-PLAN: KN-V5-C30-DERIVATION-004
0600: \SLKnowledgeBridge{逐分量求和，分清该证明与不可约性的关系。}{转移矩阵 $A$ 与分布 $\rho$，满足 $\rho_i a_{ij}=\rho_j a_{ji}$。}{$A$ 每行和为 $1$；细致平衡对所有 $i,j$ 成立。}{先固定目标状态 $j$，写 $(\rho A)_j=\sum_i\rho_i a_{ij}$。}
0601: 
0602: \begin{derivationgoalbox}
0603: 从细致平衡逐分量证明$\rho$是平稳分布，同时指出证明没有使用不可约性。
0604: \end{derivationgoalbox}
0605: \begin{derivationprepbox}
0606: 固定目标状态$j$，对所有来源状态$i$求和；再使用$A$第$j$行的和为$1$。
0607: \end{derivationprepbox}
0608: \textbf{推导第1步：写出流入量。}
0609: \[
0610: (\rho A)_j=\sum_i\rho_i a_{ij}.
0611: \]
0612: \textbf{推导第2步：逐边反向。}由细致平衡，
0613: \[
0614: \sum_i\rho_i a_{ij}=\sum_i\rho_j a_{ji}
0615: =\rho_j\sum_i a_{ji}.
0616: \]
0617: \textbf{推导第3步：使用行和。}这里$i$遍历的是矩阵第$j$行的目的状态，所以$\sum_i a_{ji}=1$，于是$(\rho A)_j=\rho_j$。对每个$j$成立，故$\rho A=\rho$。
0618: 
0619: \begin{warningbox}
0620: 细致平衡只保证平稳性。取$A=I_2$和任意严格正概率行向量$\rho=(a,1-a)$，其中$0<a<1$，所有等式$\rho_i a_{ij}=\rho_j a_{ji}$都成立，但两个状态互不可达，链可约，且平稳分布不唯一。再取$f(1)=0,f(2)=1$：从状态$1$出发时路径恒留在$1$，时间平均恒为$0$，而指定分布下的空间平均为$\sum_i\rho_i f(i)=1-a>0$。因此仅有可逆性也不保证无条件时间平均结论、非周期性或从任意初值逐步收敛；不可约性等条件必须另行验证。
0638: \phantomsection\label{struct:V5-C01-CH17}
0639: \textbf{图形辅助。}图\ref{fig:V5-C01-chain-properties}并列展示可约非周期、不可约周期和不可约非周期三种状态图；是否存在自环、图是否连通以及返回路径长度必须分别检查。
0640: % v2.3.1 figure UID: FIG-P556-02
% Proposed post-v2.3.1 polish for FIG-P556-02: protect key-label size and stroke hierarchy.
\tikzset{
  slfig-FIG-P556-02/.style={font=\fontsize{9.2pt}{11.0pt}\selectfont,every path/.append style={line cap=round,line join=round}},
  slfig-FIG-P556-02-title/.style={font=\fontsize{10.4pt}{12.4pt}\selectfont\bfseries,anchor=north},
  slfig-FIG-P556-02-formula/.style={font=\fontsize{9.2pt}{11.0pt}\selectfont},
  slfig-FIG-P556-02-answer/.style={font=\fontsize{8.8pt}{10.5pt}\selectfont,text=SLGray,anchor=south},
}
\begin{figure}[htbp]
\centering
\begin{tikzpicture}[slfig-FIG-P556-02,>=Stealth,
  every node/.style={font=\fontsize{9.2pt}{11.0pt}\selectfont,align=center},
  card/.style={draw=SLRuleGray,rounded corners=2pt,fill=white,
    minimum width=41mm,minimum height=47mm,inner sep=5.5pt,line width=.55pt},
  state/.style={circle,draw=SLBlue,fill=SLBlue!7,minimum size=6.5mm,inner sep=0pt,
    font=\fontsize{9.4pt}{11.2pt}\selectfont,line width=.76pt},
  edge/.style={-{Stealth[length=1.5mm]},draw=SLTeal,line width=.82pt},
]
  \matrix (cards) [matrix of nodes,nodes={card},column sep=3.6mm] {
    |[name=irr]| {} & |[name=per]| {} & |[name=rec]| {} \\
  };

  \node[slfig-FIG-P556-02-title,text=SLBlue] at ([yshift=-3mm]irr.north) {不可约性};
  \node[state] (i1) at ([xshift=-9mm,yshift=7mm]irr.center) {$1$};
  \node[state] (i2) at ([xshift=9mm,yshift=7mm]irr.center) {$2$};
  \draw[edge] (i1) to[bend left=18] (i2);
  \draw[edge] (i2) to[bend left=18] (i1);
  \node[slfig-FIG-P556-02-formula] at ([yshift=-7mm]irr.center) {$i\leftrightarrow j$（支撑图连通）};
  \node[slfig-FIG-P556-02-answer] at ([yshift=3mm]irr.south)
    {回答：状态能否互相到达？};

  \node[slfig-FIG-P556-02-title,text=SLTeal] at ([yshift=-3mm]per.north) {周期性};
  \draw[draw=SLTextGray,line width=.60pt] ([xshift=-14mm,yshift=7mm]per.center)--([xshift=14mm,yshift=7mm]per.center);
  \foreach \x/\lab in {-12/2,-3/4,6/6,15/8}{
    \fill[SLTeal] ([xshift=\x mm,yshift=7mm]per.center) circle (1.4pt);
  }
  \node[slfig-FIG-P556-02-formula] at ([yshift=-7mm]per.center)
    {\shortstack{$d(i)=\gcd\{n\ge1:$\\[-.4mm]$P_{ii}^{(n)}>0\}$}};
  \node[slfig-FIG-P556-02-answer] at ([yshift=3mm]per.south)
    {回答：允许哪些回返步长？};

  \node[slfig-FIG-P556-02-title,text=SLGold] at ([yshift=-3mm]rec.north) {正常返};
  \node[state,draw=SLGold,fill=SLGold!8] (ri) at ([yshift=7mm]rec.center) {$i$};
  \draw[-{Stealth[length=1.6mm]},draw=SLGold,line width=.88pt]
    (ri) to[loop right,min distance=12mm] (ri);
  \node[slfig-FIG-P556-02-formula] at ([yshift=-7mm]rec.center) {$\mathbb E_i[\tau_i^+]<\infty$};
  \node[slfig-FIG-P556-02-answer] at ([yshift=3mm]rec.south)
    {回答：平均多久回到自身？};

  \node[draw=SLRuleGray,rounded corners=2pt,fill=SLSoftGray,
    minimum width=112mm,minimum height=7mm,inner xsep=4pt,line width=.58pt,
    font=\fontsize{9.2pt}{11.0pt}\selectfont] at ([yshift=-8mm]cards.south)
    {结构性质不可相互替代：连通性、回返节律与期望回返时间必须分别检查};
\end{tikzpicture}
\caption{三类结构必须分别判断：支撑图连通性决定不可约性，正回返时长的最大公约数决定周期性}\label{fig:V5-C01-chain-properties}
\end{figure}

0641: 
0642: 图\ref{fig:V5-C01-detailed-balance}把相反方向的平稳概率流画在同一条边上，并同时给出$A=I_2$的断开反例。流量对称不等于支撑图连通。
0643: % v2.3.1 figure UID: FIG-P556-03
% Proposed post-v2.3.1 polish for FIG-P556-03: protect key-label size and stroke hierarchy.
\tikzset{
  slfig-FIG-P556-03/.style={font=\fontsize{9.2pt}{11.0pt}\selectfont,every path/.append style={line cap=round,line join=round}},
  slfig-FIG-P556-03-title/.style={font=\fontsize{10.4pt}{12.4pt}\selectfont\bfseries,text=SLBlue},
  slfig-FIG-P556-03-verdict-text/.style={font=\fontsize{9.2pt}{11.0pt}\selectfont,anchor=west},
}
\begin{figure}[htbp]
\centering
\begin{tikzpicture}[slfig-FIG-P556-03,>=Stealth,
  every node/.style={font=\fontsize{9.2pt}{11.0pt}\selectfont,align=center},
  state/.style={circle,draw=SLBlue,fill=SLBlue!7,minimum size=9mm,inner sep=0pt,
    font=\fontsize{9.4pt}{11.2pt}\selectfont,line width=.82pt},
  flow/.style={-{Stealth[length=2mm]},draw=SLTeal,line width=1pt},
  lab/.style={font=\fontsize{9.2pt}{11.0pt}\selectfont,fill=white,inner sep=1.2pt},
  verdict/.style={draw=SLRuleGray,rounded corners=2pt,fill=white,
    minimum width=38mm,minimum height=14mm,line width=.56pt},
]
  \node[slfig-FIG-P556-03-title] at (0,2.15) {逐边双向概率流};
  \node[state] (x) at (-2.15,.85) {$x$};
  \node[state] (y) at (2.15,.85) {$y$};
  \draw[flow] (x) to[bend left=18]
    node[lab,above] {$\pi(x)K(x,y)$} (y);
  \draw[flow] (y) to[bend left=18]
    node[lab,below,yshift=2.4mm] {$\pi(y)K(y,x)$} (x);
  \node[draw=SLGold,rounded corners=2pt,fill=SLGold!8,
    minimum width=63mm,minimum height=8mm,line width=.72pt,
    font=\fontsize{9.4pt}{11.2pt}\selectfont] at (0,-.25)
    {$\pi(x)K(x,y)=\pi(y)K(y,x)$};

  \node[verdict] (v1) at (-4.45,-1.65) {};
  \node[verdict] (v2) at (0,-1.65) {};
  \node[verdict] (v3) at (4.45,-1.65) {};
  \node[circle,draw=SLTeal,text=SLTeal,minimum size=5.5mm,inner sep=0pt,
    line width=.72pt,font=\fontsize{9.2pt}{11.0pt}\selectfont]
    at ([xshift=5mm]v1.west) {$\checkmark$};
  \node[slfig-FIG-P556-03-verdict-text] at ([xshift=10mm]v1.west) {可推出平稳性};
  \node[circle,draw=SLRed,text=SLRed,minimum size=5.5mm,inner sep=0pt,
    line width=.72pt,font=\fontsize{9.2pt}{11.0pt}\selectfont]
    at ([xshift=5mm]v2.west) {$\times$};
  \node[slfig-FIG-P556-03-verdict-text] at ([xshift=10mm]v2.west) {不能推出连通};
  \node[circle,draw=SLRed,text=SLRed,minimum size=5.5mm,inner sep=0pt,
    line width=.72pt,font=\fontsize{9.2pt}{11.0pt}\selectfont]
    at ([xshift=5mm]v3.west) {$\times$};
  \node[slfig-FIG-P556-03-verdict-text] at ([xshift=10mm]v3.west) {不能推出唯一};

  % 极小断开可逆链反例：A=I_2。
  \node[state,minimum size=6.5mm,font=\fontsize{9.2pt}{11.0pt}\selectfont] (c1) at (-.65,-2.80) {$1$};
  \node[state,minimum size=6.5mm,font=\fontsize{9.2pt}{11.0pt}\selectfont] (c2) at (.65,-2.80) {$2$};
  \draw[flow,line width=.68pt] (c1) to[loop left,min distance=8mm] (c1);
  \draw[flow,line width=.68pt] (c2) to[loop right,min distance=8mm] (c2);
  \node[anchor=north west,yshift=-.8mm,
    font=\fontsize{8.8pt}{10.5pt}\selectfont,text=SLGray] at (1.15,-2.80)
    {$A=I_2$：可逆但断开，平稳分布不唯一};
\end{tikzpicture}
\caption{流量对称可证平稳，却不能证连通或唯一}\label{fig:V5-C01-detailed-balance}
\end{figure}

0644: 
1024: 命题不成立。最小反例取状态空间$S=\{1,2\}$、$A=I_2$和$\rho=(a,1-a)$，其中$0<a<1$；$\rho$已归一化。对任意$i,j\in S$，
1025: \[
1026: \begin{aligned}
1027: \rho_i a_{ij}&=\rho_i\boldsymbol1\{i=j\}
1028: =\rho_j\boldsymbol1\{j=i\}=\rho_j a_{ji},\\
1029: \rho A&=\rho,\qquad
1030: \{1\},\{2\}\text{均为闭通信类}.
